\addvspace{1cm}
\subsection{Anwedung: Lösen von Differentialgleichungen}

Was sind Differentialgleichungen? Differentialgleichungen sind solche Gleichungen, die aus Ableitungen derselben Funktion mit verschieden hohem Rang bestehen. Gegeben sein müssen dafür die Anfangsbedingungen bis zur Ableitung vor der mit dem höchsten Rang. Veranschaulichen lässt sich das an einem allgemeinen Beispiel einer DGL mit der zweiten Ableitung als höchster Ableitung:

\begin{equation}
y''(t) + a y'(t) + b y(t) = f(t), \quad y(0) = y_0, \quad y'(0) = y_1
\label{eq:dgl_grundform}
\end{equation}

Um die Laplace Transformierte der Gleichung zu erhalten, wenden wir nun ein paar der oben gelernten Rechenregelen auf diese allgemein gehlatene Funktion an. Nach Anwenden der Ableitungs- (15) und Linearitätsregel (9) erhält man die Laplace-Transformierte. Auf der rechten Seite muss entsprechend die Korrespondenztabelle angewendet werden:

\begin{align*}
& \left( s^2 Y(s) - s y_0 - y_1 \right)
+ a \left( s Y(s) - y_0 \right)
+ b Y(s) = F(s) \\
\Leftrightarrow \quad
& \left( s^2 + a s + b \right) Y(s) = F(s) + (s + a) y_0 + y_1 \\
\Leftrightarrow \quad
& Y(s) = \frac{F(s)}{s^2 + a s + b} + \frac{(s + a)y_0 + y_1}{s^2 + a s + b}
\end{align*}
\cite[S.~886]{buch1}

Durch Rücktransformation erhält man dann eine Lösung des Ursprungsproblems. Wenn man im Laplace-Bereich zwei Funktionen multipliziert, entspricht das im Zeitbereich einer Faltung dieser Funktionen. Man wendet dabei das sogenannte Faltungstheorem an:

\[
\mathcal{L}^{-1} \{ F(s) \cdot G(s) \}(t) = f(t) * g(t)
\]

So ergibt sich die Lösung der Rücktransformation:

\begin{align*}
y(t) &= [\mathcal{L}^{-1}Y](t) = \left[ \mathcal{L}^{-1} \left( F(s) \cdot \frac{1}{s^2 + as + b} \right) \right](t) 
+ \left[ \mathcal{L}^{-1} \left( \frac{(s+a)y_0 + y_1}{s^2 + as + b} \right) \right](t) \\
&= [\mathcal{L}^{-1}F(s)](t) * \left[ \mathcal{L}^{-1} \left( \frac{1}{s^2 + as + b} \right) \right](t) 
+ \left[ \mathcal{L}^{-1} \left( \frac{(s+a)y_0 + y_1}{s^2 + as + b} \right) \right](t)
\end{align*}
\cite[S.~886]{buch1}

Dazu gegeben sei ein Zahlenbeispiel mit einfacher Ableitung und einer Exponentialfunktion im Eingang:

\begin{align}
    6\dot{y}(t) + y(t) = 43 \, \mathrm{e}^{-7t} \\
    \text{mit } y(0) = 3
\end{align}

Wir transformieren beide Seiten, zum einen mit Hilfe der Ableitungsregel und zum anderen mit Hilfe der Korrespondenztablle für die Exponentialfunktion:
\begin{itemize}
  \item $\mathcal{L}\{\dot{y}(t)\} = sY(s) - y(0)$
  \item $\mathcal{L}\{y(t)\} = Y(s)$
  \item $\mathcal{L}\{e^{-at}\} = \frac{1}{s + a}$
\end{itemize}

Dann ergibt sich:
\[
6(sY(s) - 3) + Y(s) = \frac{43}{s + 7}
\]

\textbf{Rücktransformation mit der Faltungsregel}

Wir berechnen die Rücktransformation
\[
y(t) = \mathcal{L}^{-1} \left\{ \frac{43}{(s + 7)(6s + 1)} \right\} + \mathcal{L}^{-1} \left\{ \frac{18}{6s + 1} \right\}
\]

Zunächst der zweite Term. Dafür bedienen wir uns der Korrespondenztabelle:
\[
\frac{18}{6s + 1} = 3 \cdot \frac{1}{s + \frac{1}{6}} \quad \Rightarrow \quad \mathcal{L}^{-1} \left\{ \frac{18}{6s + 1} \right\} = 3 \cdot e^{-\frac{1}{6}t}
\]

Nun wenden wir die Faltungsregel auf den ersten Term an:

\[
\mathcal{L}^{-1} \left\{ \frac{43}{(s + 7)(6s + 1)} \right\}
= 43 \cdot \mathcal{L}^{-1} \left\{ \frac{1}{s + 7} \cdot \frac{1}{6s + 1} \right\}
\]

\[
= 43 \cdot \left( e^{-7t} * \frac{1}{6} e^{-\frac{1}{6}t} \right)
= \frac{43}{6} \cdot \int_0^t e^{-7\tau} \cdot e^{-\frac{1}{6}(t - \tau)} \, d\tau
\]

\[
= \frac{43}{6} \cdot e^{-\frac{1}{6}t} \cdot \int_0^t e^{-\frac{41}{6}\tau} \, d\tau
= \frac{43}{6} \cdot e^{-\frac{1}{6}t} \cdot \left[ \frac{-1}{\frac{41}{6}} e^{-\frac{41}{6}\tau} \right]_0^t
\]

\[
= \frac{43}{6} \cdot e^{-\frac{1}{6}t} \cdot \frac{6}{41} \left( 1 - e^{-\frac{41}{6}t} \right)
= \frac{43}{41} \cdot e^{-\frac{1}{6}t} \cdot \left( 1 - e^{-\frac{41}{6}t} \right)
\]

Zusammengefügt bekommen wir schließlich unsere allgemeine Lösung:

\[
\Rightarrow y(t) = \frac{43}{41} \cdot e^{-\frac{1}{6}t} \cdot \left( 1 - e^{-\frac{41}{6}t} \right) + 3 \cdot e^{-\frac{1}{6}t}
\]
    
\addvspace{1cm}
\textbf{Unterschied in der Anwendung der Ableitungsregel auf die Fouriertransformation vs die Laplace Transformation}

Die Fourier-Transformation unterscheidet sich in Bezug auf die Ableitungsregel von der Laplace-Transformation, da bei ihr keine Anfangsbedingungen berücksichtigt werden – was darauf zurückzuführen ist, dass über den gesamten Zahlenraum von $-\infty$ bis $\infty$ integriert wird, statt wie bei der Laplace-Transformation ab null. 

Außerdem sind Funktionen mit exponentiellem Wachstum zwar Laplace-, aber nicht Fourier-transformierbar. Dadurch lassen sich bestimmte Lösungen mit der Fourier-Transformation nicht erfassen. 

Da exponentielle Funktionen jedoch häufig als Grundbausteine in Lösungen von Differenzialgleichungen auftreten, bietet die Laplace-Transformation hier einen klaren Vorteil \cite[891]{buch1}.
